%% ==========================================================================
%% abnt.tex — Padrão visual ABNT (NBR 14724) do Projeto Integrador
%% Parte da extensão Quarto pi-abnt. Não edite: atualize com `quarto update extension`.
%% ==========================================================================

%% --- Parágrafos: recuo de 1,25 cm, inclusive no primeiro parágrafo ---------
\usepackage{indentfirst}
\setlength{\parindent}{1.25cm}
\setlength{\parskip}{0pt}

%% --- Legendas: "Figura 1 – Título" acima, fonte menor ---------------------
\usepackage{caption}
\captionsetup{labelsep=endash, font=small, justification=centering,
              singlelinecheck=false, skip=4pt}

%% --- Numeração sequencial de figuras/tabelas (não por capítulo) ------------
\counterwithout{figure}{chapter}
\counterwithout{table}{chapter}

%% --- Idioma: português do Brasil (hifenização, "Figura", "Tabela" etc.) ---
%% Usa os arquivos .ini do próprio babel: não depende do pacote babel-portuges.
\usepackage[provide=*]{babel}
\babelprovide[import, main]{brazilian}

%% --- Títulos das seções (NBR 6024) -----------------------------------------
%% 1 PRIMÁRIA (negrito, maiúsculas) · 1.1 SECUNDÁRIA (maiúsculas)
%% 1.1.1 Terciária (negrito) · 1.1.1.1 Quaternária (itálico)
\makeatletter
\renewcommand{\@makechapterhead}[1]{%
  {\normalfont\normalsize\bfseries\noindent
   \ifnum\c@secnumdepth>\m@ne\thechapter\hspace{0.6em}\fi
   \MakeUppercase{#1}\par\vspace{1.5em}}}
\renewcommand{\@makeschapterhead}[1]{%
  {\normalfont\normalsize\bfseries\centering\MakeUppercase{#1}\par\vspace{1.5em}}}
\renewcommand{\section}{\@startsection{section}{1}{0pt}%
  {1.5em}{1em}{\normalfont\normalsize\MakeUppercase}}
\renewcommand{\subsection}{\@startsection{subsection}{2}{0pt}%
  {1.5em}{1em}{\normalfont\normalsize\bfseries}}
\renewcommand{\subsubsection}{\@startsection{subsubsection}{3}{0pt}%
  {1.5em}{1em}{\normalfont\normalsize\itshape}}
\renewcommand{\@seccntformat}[1]{\csname the#1\endcsname\hspace{0.6em}}
\makeatother
\setcounter{secnumdepth}{4}
\setcounter{tocdepth}{4}

%% --- Paginação: contada desde a folha de rosto, exibida a partir da --------
%% --- Introdução, no canto superior direito ---------------------------------
\usepackage{fancyhdr}
\fancypagestyle{abntnumerada}{\fancyhf{}\fancyhead[R]{\small\thepage}%
  \renewcommand{\headrulewidth}{0pt}}
\pagestyle{empty}
\fancypagestyle{plain}{\fancyhf{}\renewcommand{\headrulewidth}{0pt}}
\newcommand{\abnttextual}{%
  \clearpage
  \pagestyle{abntnumerada}%
  \fancypagestyle{plain}{\fancyhf{}\fancyhead[R]{\small\thepage}%
    \renewcommand{\headrulewidth}{0pt}}}

%% --- Títulos sem número (Resumo, Abstract, Referências, Apêndices) ---------
\newcommand{\abntsemnumero}[2]{% #1 = título, #2 = 1 entra no sumário
  \clearpage
  \phantomsection
  \ifnum#2=1 \addcontentsline{toc}{chapter}{\MakeUppercase{#1}}\fi
  \begin{center}\textbf{\MakeUppercase{#1}}\end{center}
  \vspace{1em}}

%% --- Listas e sumário com título centralizado ------------------------------
\AtBeginDocument{%
  \renewcommand{\contentsname}{Sumário}%
  \renewcommand{\listfigurename}{Lista de figuras}%
  \renewcommand{\listtablename}{Lista de tabelas}%
  \renewcommand{\figurename}{Figura}%
  \renewcommand{\tablename}{Tabela}}
\newcommand{\abntlistas}{%
  \renewcommand{\contentsname}{Sumário}%
  \renewcommand{\listfigurename}{Lista de figuras}%
  \renewcommand{\listtablename}{Lista de tabelas}%
  \listoffigures\clearpage
  \listoftables\clearpage
  \tableofcontents}

%% --- "Fonte:" abaixo de figuras e tabelas ----------------------------------
\newenvironment{abntfonte}{\par\vspace{-0.6em}\begin{center}\small}{\end{center}\vspace{0.4em}}
